\section{Experimental design}
\label{sec:design}

The experiments are not a chronology. Before training anything, the observations of
Section~\ref{sec:data} were turned into an ordered list of hypotheses, each changing
exactly one variable and each falsifiable by a single number:

\begin{enumerate}\itemsep0pt
    \item a plain encoder--decoder already finds the main filaments, and gives a
          baseline worth comparing against;
    \item with fewer than 500 observations the binding constraint is the encoder,
          so ImageNet features should beat random initialisation;
    \item Dice is blind to fragmentation, so a loss that is not -- Tversky plus
          clDice -- should raise PQ where Dice cannot;
    \item part of what remains is the H-alpha \emph{domain} rather than the loss;
    \item what limits the score is \emph{scale}: at $512 \times 512$ a filament six
          native pixels wide is one and a half pixels wide.
\end{enumerate}

Hypotheses 3 and 5 make opposite predictions about where the residual error lives,
and only one of them survives contact with the data. Later hypotheses were derived
the same way, each from the error analysis of the run before it.

Two rules kept the study honest. A candidate is compared against the previous best
only \emph{at its own operating point}, re-searched from scratch, because a
different loss or resolution shifts the calibration and reusing an old threshold
would falsify a good idea for the wrong reason. And a gain is believed only if it
survives \textbf{cross-selection}: the operating point is chosen on one half of the
validation files and scored on the other. Measured that way, on 169 annotator
entries, \textbf{any gain below $\approx 0.005$ is indistinguishable from selection
noise} -- a threshold that later killed three mechanisms whose aggregate numbers
had looked positive.

Before starting, we also wrote down what we were going to explore and how long it
would take, so that the study would fit in the time we had
(Table~\ref{tab:plan}).

\begin{table}[ht]
    \centering
    \caption{The plan: what we decided to explore, and the cost of one run.}
    \label{tab:plan}
    \small\begin{tabular}{lll}
        \toprule
        What we explore & Values considered & Cost \\
        \midrule
        input resolution      & 512, 1024, 1536, native 2048 & 10--35 min / run \\
        encoder               & U-Net from scratch, ResNet18, ResNet34 & 10--35 min / run \\
        loss function         & BCE+Dice, +Tversky, +clDice, instance Dice, \texttt{pos\_weight} & 15--30 min / run \\
        augmentation          & D4 symmetries, photometric jitter, scale jitter & included above \\
        training paradigm     & pixel-wise U-Net vs.\ YOLOv8 detector & 30 min -- 2.5 h / run \\
        post-processing       & threshold, closing, min.\ size, confidence gate & 10 min, no training \\
        \bottomrule
    \end{tabular}
\end{table}

